а) У поля $F$ конечной характеристики $char F$ -- простое число.

Предположим, что $char F = ab$, где $a, b > 1$. Обозначим
$\overline{x} = \underbrace{1 + 1 + \dots + 1}_{x раз}$.

Тогда получается, что $\overline{a} \overline{b} = 0$.
Ни одно из них не равно нулю, так как $a < char F$ и $b < char F$.
Получается, что $\overline{a}$ и $\overline{b}$ -- делители нуля, что
противоречит тому, что $F$ -- поле.

б) Если существует нетривиальный гомоморфизм полей 
$\varphi: F \to K$, то $char F = char K$.

Поскольку $\varphi(1) = 1$, то 
$\forall a \ \ \varphi(\overline{a}_F) = \overline{a}_K$,
поэтому характеристики полей $F$ и $K$ совпадают. 

в) Любое конечное поле имеет положительную характеристику.

Предположим, что $char F = 0$, $|F| < \infty$.

Так как поле конечно, то среди последовательности 
$\overline{1}, \overline{2}, \overline{3}, \dots$ найдутся два
одинаковых элемента. Пусть $\overline{k} = \overline{l}$, $l > k$.
Но тогда $\overline{l - k} = 0$. Противоречие с тем, что $char F = 0$.

г) Существует ли бесконечное поле положительной характеристики?

Да, например $Quot Z_2[x] \cong \Z_2(x)$. $char \Z_2(x) = 2$

д) (Нетривиальный) гомоморфизм полей $\varphi: F \to K$ является инъективным

$Ker \varphi$ является идеалом $F$. В полях возможны только тривиальные идеалы:
$\{0\}$ или $F$. Если $Ker \varphi = F$, то гомоморфизм тривиальный.
Получается, что $Ker \varphi = \{0\}$, то есть $\varphi$ -- инъективный.

е) Если $L \supset F$ -- расширение полей, 
то $L$ является линейным пространством над $F$.

По определению поля, $(L, +)$ -- абелева группа, поэтому проверим
аксиомы только для умножения на элементы из $F$:

1) $f \cdot (l_1 + l_2) = f l_1 + f l_2$, 
$(f_1 + f_2) l = f_1 l + f_2 l$

2) $1_F \cdot l = 1$

3) $f_1 (f_2 l) = (f_1 f_2) l$
